% SPDX-FileCopyrightText: Copyright © 2026 Idiap Research Institute <contact@idiap.ch>
% SPDX-FileContributor: Andre Freitas <andre.freitas@idiap.ch>
% SPDX-License-Identifier: MIT
% Neuro-symbolic AI Group
% File Documentation
%
% Motivation:
%   Protect Mathformer behavior while the framework evolves across Lean, ATP, KB, SymPy, docs, and skills.
%
% Intent:
%   Encode regression checks that catch parser breakage, semantic drift, and missing workflow evidence.
%
% Function:
%   Contains executable tests, fixtures, and manifests used for local and CI validation.
%
% Maintenance Notes:
%   Keep this header synchronized with major responsibility changes, and prefer updating tests or examples when behavior changes.

% mathformer-case-id: wiki_sum_first_odd_numbers
% expected-lean-name: WikiLevelGold.wiki_sum_first_odd_numbers
% controlled-suite: wiki_level_latex_to_lean
% difficulty: medium
\documentclass[11pt]{article}
\usepackage[T1]{fontenc}
\usepackage[utf8]{inputenc}
\usepackage{lmodern}
\usepackage{amsmath,amssymb,amsthm}

\newtheorem{theorem}{Theorem}

\title{The Sum of the First Odd Numbers}
\author{}
\date{}

\begin{document}

\maketitle

\section*{Controlled Origin Annotation}
This is a synthetic controlled case written in a Wikipedia-style exposition.
It is intended to test indexed sums, natural-number induction, arithmetic
normalization, and the convention that the first $n$ odd numbers are
$1,3,\ldots,2n-1$.

\section*{Statement and Proof}

\begin{theorem}[Sum of the first odd numbers]
For every natural number $n$,
\[
  1 + 3 + 5 + \cdots + (2n-1) = n^2.
\]
\end{theorem}

\begin{proof}
We prove the result by induction on $n$. For $n=0$, the sum has no terms and
is $0$, which is $0^2$.

Assume that the sum of the first $n$ odd numbers is $n^2$. The next odd number
is $2n+1$, so the sum of the first $n+1$ odd numbers is
\[
  n^2 + (2n+1).
\]
But
\[
  n^2 + 2n + 1 = (n+1)^2.
\]
Therefore the result holds for $n+1$. By induction, it holds for every natural
number $n$.
\end{proof}

\section*{Lean Target}
Formalize the first $n$ odd numbers as the finite sum
\[
  \sum_{i=0}^{n-1} (2i+1).
\]
The expected Lean theorem is \texttt{WikiLevelGold.wiki\_sum\_first\_odd\_numbers}.

\end{document}
